\begin{frame}[t]{ROM --- Read-Only Memory}
  \begin{block}{Características}
    Memória não volátil destinada principalmente à leitura. Mantém informações após o desligamento.
  \end{block}
  \begin{itemize}
    \item Pode guardar firmware de inicialização, tabelas fixas e identificações.
    \item Em ROM clássica, o conteúdo é definido na fabricação.
    \item O termo ROM também é usado para memória de programa não volátil.
  \end{itemize}
  \begin{exampleblock}{Exemplo}
    Um \textit{bootloader} pode estar armazenado em uma região de memória somente leitura do dispositivo.
  \end{exampleblock}
\end{frame}
